% ============================================================
%  SCHOLA DOMESTICA — Level 1, Week 2, Lesson 10
%  Topic: Numbers 0–10 Review & Number Line Introduction
%  10–12 problems | No speed drill | Friday consolidation
%  Compile with XeLaTeX
% ============================================================
\documentclass[12pt,letterpaper]{article}
\usepackage{../sd_style}

\renewcommand{\lessonnum}{10}
\renewcommand{\lessonweek}{2}

\begin{document}

\lessonheader{2}{10}{Numbers 0–10: Review \& the Number Line}

\begin{instructbox}
\noindent\textbf{\color{sd-blue}Instruction} \textit{(read aloud):}\\[4pt]
Wonderful work this week!
Today we review all the numbers from 0 to 10 and meet the \textbf{number line}.
A number line shows numbers in order from left to right —
the farther right, the greater the number.
\end{instructbox}

\vspace{8pt}
% Number line
\noindent\textbf{\color{sd-blue}The Number Line:}
\vspace{8pt}

\begin{center}
\begin{tikzpicture}
  \draw[sd-blue, very thick, ->] (-0.3,0) -- (11.5,0);
  \foreach \n in {0,...,10}{
    \draw[sd-blue, thick] (\n, -0.15) -- (\n, 0.15);
    \node[below] at (\n, -0.2) {\textbf{\n}};
  }
\end{tikzpicture}
\end{center}

\goldrule
\noindent\textbf{\color{sd-blue}Practice}
\vspace{6pt}

\prob Count and write each number.\\[8pt]
\begin{tabular}{ccccc}
\IfFileExists{stars_6.jpg}{\includegraphics[height=1.2cm]{stars_6.jpg}}{%
  \onestar\onestar\onestar\onestar\onestar\onestar} &
\IfFileExists{stars_8.jpg}{
\includegraphics[height=1.2cm]{stars_8.jpg}}{%
  \onestar\onestar\onestar\onestar\onestar\onestar\onestar\onestar} &
\IfFileExists{stars_10.jpg}{
\includegraphics[height=1.2cm]{stars_10.jpg}}{%
  \onestar\onestar\onestar\onestar\onestar\onestar\onestar\onestar\onestar\onestar} &
\IfFileExists{stars_7.jpg}{\includegraphics[height=1.2cm]{stars_7.jpg}}{%
  \onestar\onestar\onestar\onestar\onestar\onestar\onestar} &
\IfFileExists{stars_9.jpg}{
\includegraphics[height=1.2cm]{stars_9.jpg}}{%
  \onestar\onestar\onestar\onestar\onestar\onestar\onestar\onestar\onestar}\\[4pt]
\ansblank & \ansblank & \ansblank & \ansblank & \ansblank
\end{tabular}

\vspace{12pt}
\prob Write the missing numbers on the number line below.\\[8pt]
\begin{center}
\begin{tikzpicture}
  \draw[sd-blue, very thick, ->] (-0.3,0) -- (11.5,0);
  \foreach \n in {0,...,10}{
    \draw[sd-blue, thick] (\n, -0.15) -- (\n, 0.15);
  }
  % Show some, hide others
  \foreach \n in {0,2,4,6,8,10}{
    \node[below] at (\n, -0.25) {\textbf{\n}};
  }
  \foreach \n in {1,3,5,7,9}{
    \node[below] at (\n, -0.25) {\underline{\hspace{0.5cm}}};
  }
\end{tikzpicture}
\end{center}

\vspace{12pt}
\prob What number is \textbf{one more} than 7? \quad\ansblank

\vspace{10pt}
\prob What number is \textbf{one less} than 9? \quad\ansblank

\vspace{10pt}
\prob What number is \textbf{between} 5 and 7? \quad\ansblank

\vspace{10pt}
\prob Put these numbers in order, greatest to least:\\[6pt]
\hspace{2cm}{\Large 4 \quad 9 \quad 2 \quad 7 \quad 6}\\[6pt]
\hspace{2cm}\ansblank\quad\ansblank\quad\ansblank\quad\ansblank\quad\ansblank

\vspace{10pt}
\prob \textit{(Teacher reads)} A little mouse ran to the right on the number line
starting at 3 and stopped at 8.
How many hops did she take? \quad\ansblank

\vspace{10pt}
\prob Circle every number that is \textbf{greater than 6}:\\[6pt]
\hspace{2cm}{\Large 0 \quad 5 \quad 7 \quad 10 \quad 3 \quad 9 \quad 6 \quad 8}

\vspace{10pt}
\prob Circle every number that is \textbf{less than 4}:\\[6pt]
\hspace{2cm}{\Large 5 \quad 1 \quad 3 \quad 8 \quad 0 \quad 4 \quad 2 \quad 7}

\vspace{10pt}
\prob Write the numbers 0 through 10 in your neatest hand:\\[8pt]
\noindent\fbox{\rule{0pt}{1.8cm}\hspace{12.5cm}}

\goldrule

\begin{checkbox}
\noindent\textbf{\color{sd-green}Knowledge Check}\\[4pt]
\setcounter{probnum}{0}
\prob Count to 10 and back to 0 aloud. \quad $\square$~Done\\[8pt]
\prob What number comes between 8 and 10? \quad\ansblank\\[8pt]
\prob Which is greater: 6 or 9? \quad\ansblank
\end{checkbox}

\vfill
\noindent{\color{sd-blue}$\bigstar$ \textbf{Enrichment — Roman Numerals I–X:}}\\
You now know all ten Roman numerals!\\[4pt]
\begin{center}
\begin{tabular}{cccccccccc}
\textbf{I}&\textbf{II}&\textbf{III}&\textbf{IV}&\textbf{V}&
\textbf{VI}&\textbf{VII}&\textbf{VIII}&\textbf{IX}&\textbf{X}\\[2pt]
1&2&3&4&5&6&7&8&9&10
\end{tabular}
\end{center}
\vspace{4pt}
Write your age in Roman numerals: \quad\ansblank

\vspace{4pt}\hrule\vspace{3pt}
{\footnotesize\color{gray}\textit{Teacher note: Week 2 mastery target — student reads,
writes, and orders all numbers 0–10 without hesitation and understands
the number line as a left-to-right sequence of increasing values.
If not solid, do not advance to Week 3 before reviewing.}}

\end{document}
